% Part: turing-machines
% Chapter: machines-computations
% Section: introduction

\documentclass[../../../include/open-logic-section]{subfiles}

\begin{document}

\olfileid{tur}{mac}{int}
\olsection{परिचय}

$\Nat$ पासून $\Nat$ मध्ये जाणारे एखादे फलन \emph{संगणनक्षम} असणे म्हणजे
नेमके काय? या प्रश्नाला सर्वप्रथम दिलेल्या आणि सर्वाधिक परिचित
उत्तरांपैकी एक असे आहे: एखादे फलन ट्यूरिंग यंत्राने काढता येत असेल, तर ते
संगणनक्षम असते. ॲलन ट्यूरिंग यांनी 1936 मध्ये ही संकल्पना मांडली.
ट्यूरिंग यंत्र हे \emph{संगणनाचे प्रतिमान} आहे---म्हणजे ``संगणनाची
कार्यपद्धती'' ही कल्पना गणिती नेमकेपणाने परिभाषित करण्याचा एक मार्ग.
त्या कल्पनेचा नेमका अर्थ काय, यावर मतभेद असले तरी ट्यूरिंग यंत्रे ही
संगणनाच्या कार्यपद्धती निर्दिष्ट करण्याची एक पद्धत आहेत, यावर व्यापक
सहमती आहे. ``ट्यूरिंग यंत्र'' या नावामुळे हलणारे भाग असलेल्या भौतिक
यंत्राची प्रतिमा डोळ्यांसमोर येऊ शकते; पण काटेकोरपणे पाहता ट्यूरिंग
यंत्र ही पूर्णपणे गणिती रचना आहे. म्हणून ती संगणनाच्या कार्यपद्धतीची
आदर्शीकृत मांडणी करते. उदाहरणार्थ, ट्यूरिंग यंत्राला लागणाऱ्या वेळेवर
किंवा स्मृतीवर आपण कोणतेही बंधन घालत नाही: एखाद्या संगणनाला विश्वातील
अणूंपेक्षा जास्त साठवण-जागा किंवा पायऱ्या लागणार असल्या तरी ट्यूरिंग
यंत्र ते संगणन करू शकते.

\begin{explain}
ट्यूरिंग यंत्राचा विचार एका विशिष्ट प्रकारच्या काल्पनिक यंत्रणेसाठीचा
कार्यक्रम म्हणून करणे बहुधा सर्वांत सोयीचे ठरेल. या यंत्रणेत एक
\emph{फीत} आणि एक \emph{वाचन-लेखन शीर्ष} असते. आपल्या ट्यूरिंग
यंत्रांच्या आवृत्तीत फीत एका दिशेने (उजवीकडे) अनंत असते आणि ती
\emph{घरांमध्ये} विभागलेली असते. प्रत्येक घरात सांत \emph{वर्णमालेतील}
एखादे चिन्ह असू शकते. अशा वर्णमालेत कितीही भिन्न चिन्हे असू शकतात;
पण आपल्याला मुख्यतः तीनच पुरेशी पडतील: $\TMendtape$, $\TMblank$ आणि
$\TMstroke$. यंत्रणा सुरू करताना फीत रिकामी असते (म्हणजे प्रत्येक घरात
$\TMblank$ हे चिन्ह असते); अपवाद म्हणजे सर्वांत डावीकडील घर, ज्यात
$\TMendtape$ असते, आणि \emph{आदान} असलेली सांत संख्येने घरे. कोणत्याही
क्षणी यंत्रणा सांत संख्येतील एखाद्या \emph{अवस्थेत} असते. सुरुवातीला
शीर्ष सर्वांत डावीकडील घर वाचत असते आणि यंत्रणा निर्दिष्ट
\emph{प्रारंभिक अवस्थेत} असते. यंत्रणेच्या प्रत्येक पायरीवर शीर्ष
सध्या वाचत असलेल्या घरातील चिन्ह, यंत्रणेची अवस्था आणि ट्यूरिंग
यंत्राचा कार्यक्रम मिळून पुढे काय होईल ते ठरवतात. हा कार्यक्रम एका
आंशिक फलनाने दिलेला असतो: त्याला अवस्था~$q$ आणि चिन्ह~$\sigma$ ही
आदाने दिल्यावर ते $\tuple{q', \sigma',
  D}$ ही त्रयी देते. यंत्रणा $q$ अवस्थेत असताना $\sigma$ चिन्ह
वाचत असेल, तर ती सध्याच्या घरातील चिन्हाच्या जागी $\sigma'$ लिहिते;
$D$ हे $\TMleft$, $\TMright$ किंवा $\TMstay$ यांपैकी कोणते आहे,
त्यानुसार शीर्ष डावीकडे, उजवीकडे किंवा त्याच जागी राहते; आणि यंत्रणा
$q'$ अवस्थेत जाते.

उदाहरणार्थ, \olref{fig:tm} मधील स्थिती पाहा.
\begin{figure}
  \olasset{\olpath/assets/diagrams/turing-machine.tikz}
  \caption{कार्यक्रम चालवत असलेले ट्यूरिंग यंत्र.}
  \ollabel{fig:tm}
\end{figure}
ट्यूरिंग यंत्राच्या फितीच्या दिसणाऱ्या भागात सर्वांत डावीकडील घरात
फितीच्या टोकाचे $\TMendtape$ चिन्ह आहे; त्यानंतर तीन $1$, एक $0$ आणि
आणखी चार $1$ आहेत. शीर्ष डावीकडून तिसरे घर वाचत आहे; त्या घरात
$1$ आहे आणि यंत्रणा $q_1$ अवस्थेत आहे---म्हणजे ``यंत्र $q_1$
अवस्थेत $1$ वाचत आहे'' असे आपण म्हणतो. ट्यूरिंग यंत्राचा कार्यक्रम
$\tuple{q_1, 1}$ या आदानासाठी $\tuple{q_2,
0, \TMstay}$ ही त्रयी देत असेल, तर यंत्र तिसऱ्या घरातील $1$ च्या
जागी $0$ लिहील, वाचन-लेखन शीर्ष त्याच जागी ठेवील आणि $q_2$
अवस्थेत जाईल. त्यानंतर कार्यक्रम $\tuple{q_2, 0}$ या आदानासाठी
$\tuple{q_3, 0, \TMright}$ देत असेल, तर यंत्र त्या $0$ च्या
जागी पुन्हा $0$ लिहील (म्हणजे प्रत्यक्षात शीर्षाखालील फितीतील मजकूर
अपरिवर्तित राहील), एक घर उजवीकडे सरकेल आणि $q_3$ अवस्थेत जाईल.
अशाच रीतीने पुढे चालेल.

एखादी अवस्था $q_n$ आणि चिन्ह $\sigma$ अशी आढळली की
$\tuple{q_n, \sigma}$ साठी कोणतीही सूचना नाही---म्हणजे
$\tuple{q_n,\sigma}$ या आदानावर संक्रमण फलन अपरिभाषित आहे---तर यंत्र
\emph{थांबते} असे आपण म्हणतो.
दुसऱ्या शब्दांत, यंत्राला अमलात आणण्यासाठी कोणतीही सूचना उरत नाही
आणि त्या टप्प्यावर त्याचे कार्य थांबते. काही वेळा थांबणे एका विशिष्ट
थांबण्याच्या अवस्थेने~$h$ दर्शवतात. याचे अधिक तपशीलवार प्रदर्शन पुढे
केले जाईल.
\end{explain}

\begin{digress}
ट्यूरिंग यांच्या ``On computable numbers'' या शोधनिबंधाचे वैशिष्ट्य
असे की त्यात त्यांनी केवळ औपचारिक व्याख्याच दिली नाही, तर ती व्याख्या
संगणनक्षमतेची अंतर्ज्ञानी कल्पना योग्यरीत्या पकडते, असा युक्तिवादही
केला. व्याख्येवरून ट्यूरिंग यंत्राने काढता येणारे कोणतेही फलन अंतर्ज्ञानी
अर्थाने संगणनक्षम आहे, हे स्पष्ट व्हावे. ट्यूरिंग यांनी याचा उलट
दिशेचा दावा---म्हणजे आपण स्वाभाविकपणे संगणनक्षम मानू असे कोणतेही फलन
अशा यंत्राने काढता येते---यासाठी तीन प्रकारचे युक्तिवाद दिले. त्यांच्या
शब्दांत ते असे:
\begin{enumerate}
\item अंतःप्रज्ञेला थेट आवाहन.
\item दोन व्याख्या सममूल्य आहेत, याची सिद्धता (नवी व्याख्या
  अंतःप्रज्ञेला अधिक पटत असल्यास).
\item संगणनक्षम संख्यांच्या मोठ्या वर्गांची उदाहरणे देणे.
\end{enumerate}
आपला उद्देश संगणनक्षमता ``तत्त्वतः'' परिभाषित करण्याचा आहे: म्हणजे
वेळ आणि जागा यांच्या व्यावहारिक मर्यादा विचारात न घेता. अर्थात,
संगणनक्षमतेची ही सर्वांत व्यापक व्याख्या मिळाल्यानंतर मर्यादित
संसाधनांतील संगणनाचाही अभ्यास करता येतो. ``संगणकीय जटिलता'' म्हणून
ओळखल्या जाणाऱ्या विषयाचा हा मध्यवर्ती भाग आहे.
\end{digress}

\begin{history}
ॲलन ट्यूरिंग यांनी 1936 मध्ये ट्यूरिंग यंत्रांची कल्पना मांडली.
त्या काळी त्यांना प्रथम-क्रम तर्कशास्त्रातील निर्णयक्षमतेचा प्रश्न
विशेष महत्त्वाचा वाटत होता; तरी संगणकाच्या रचनेच्या पायावरील निर्णायक
शोधनिबंध म्हणून त्यांच्या त्या निबंधाचे वर्णन झाले आहे. त्या निबंधात
ट्यूरिंग यांनी संगणनक्षम वास्तव संख्यांवर---म्हणजे ज्यांचे दशांश
विस्तार संगणनक्षम आहेत अशा संख्यांवर---लक्ष केंद्रित केले. पण
नैसर्गिक संख्यांवरील संगणनक्षम फलनांसाठी आणि पुढील गोष्टींसाठीही या
संकल्पना रूपांतरित करणे कठीण नाही, असे त्यांनी नमूद केले. लक्षात घ्या:
पहिला कार्यरत सर्वसाधारण वापराचा संगणक 1941 मध्ये (जर्मन कॉनराड
झुसे यांनी आपल्या पालकांच्या दिवाणखान्यात) बांधला जाण्याच्या पूर्ण पाच
वर्षे आधी हे झाले. 1936 नंतर सात वर्षांनी (1943) ट्यूरिंग आणि ब्लेच्ली
पार्कमधील त्यांचे सहकारी यांनी संकेतभेदनासाठी कोलॉसस बांधला;
अमेरिकन ENIAC
(1945) हा नऊ वर्षांनी आला; मँचेस्टरमधील पहिले ब्रिटिश सर्वसाधारण
वापराचे यंत्र---Manchester Small-Scale Experimental Machine---बारा
वर्षांनी (1948) बांधले गेले; आणि अमेरिकन BINAC (1949) ची पहिली
चाचणी तेरा वर्षांनी झाली. Manchester SSEM हा संग्रहित-कार्यक्रम
असलेला पहिला संगणक होता---आधीच्या यंत्रांमध्ये प्रत्येक नव्या कामासाठी
तारा हाताने नव्याने जोडाव्या लागत.
\end{history}

\end{document}
